\textbf{Objetivo general.} Construir y documentar un modelo reproducible para estimar el rendimiento de las 59 parcelas objetivo del Reto AgroCebada FIRA 2026 a partir de atributos agroambientales conocidos y de 138 rendimientos observados, con énfasis en precisión predictiva, trazabilidad y prevención de filtración de información.

\textbf{Objetivos específicos.} Primero, integrar datos satelitales, climáticos, topográficos y geológicos en registros por parcela, respetando ventanas temporales, unidades y procedencia de cada sensor. Segundo, contrastar alternativas de tratamiento de valores faltantes, como mediana, media, imputación espacial, temporal y espaciotemporal, sin utilizar el rendimiento de la parcela retenida. Tercero, comparar regresiones globales regularizadas, expertos geográficos y compuertas basadas en perfiles de composición geológica, incluyendo el examen de K-means y DBSCAN. Cuarto, investigar la influencia del número de vecinos y la penalización Ridge sobre conjuntos de validación idénticos. Quinto, estudiar correcciones residuales y ajustes conservadores por análogos con controles explícitos de soporte. Sexto, reportar RMSE, MAE, sesgo, coeficiente de determinación predictivo, sensibilidad a parcelas extremas y autocorrelación espacial. Finalmente, separar las inferencias válidas sobre las 138 parcelas conocidas de los resultados todavía no verificables para las 59 parcelas sin etiqueta.
